\section{Introduction}
\label{intro:section}

Calder\'on's inverse boundary problem asks whether boundary measurements
determine the coefficients of an elliptic equation. In its geometric
form, the unknown is a Riemannian metric and the measurements are the
Dirichlet energy of harmonic extensions. A change of interior coordinates
that fixes the measured boundary points preserves these measurements.
We prove that this is the only ambiguity for smooth metrics in dimension
at least three, even when both the prescribed values and the observations
are confined to one arbitrary open boundary patch.

\subsection{The measurement and the theorem}

Let $M$ be a compact connected smooth $n$-manifold with nonempty smooth
boundary, and let $g$ be a smooth positive-definite Riemannian metric on
$M$, including its boundary. All volume and surface integrals use
densities, so no orientation is required. For $f\in H^{1/2}(\partial M)$,
let $u_f^g\in H^1(M)$ be the unique weak solution of
\[
 -\Delta_g u_f^g=0\quad\text{in }M,\qquad
 u_f^g|_{\partial M}=f.
\]
Here $\Delta_g=\operatorname{div}_g\nabla_g$. For a nonempty relatively
open set $\Gamma\subset\partial M$, put
\[
 H^{1/2}_{co}(\Gamma)
 =\{f\in H^{1/2}(\partial M):\supp f\subset\Gamma\}.
\]
Equip this space with the norm inherited from $H^{1/2}(\partial M)$.
The local energy Dirichlet-to-Neumann map is the bounded operator
\[
 \Lambda_{g,\Gamma}:H^{1/2}_{co}(\Gamma)
       \longrightarrow (H^{1/2}_{co}(\Gamma))^*
\]
defined by
\begin{equation}\label{intro:energy}
 \langle\Lambda_{g,\Gamma}f,h\rangle
 =\int_M\langle\dd u_f^g,\dd u_h^g\rangle_g\,\dd V_g.
\end{equation}
For smooth supported data, Green's formula writes this as
\begin{equation}\label{intro:density}
 \langle\Lambda_{g,\Gamma}f,h\rangle
 =\int_\Gamma h\,\partial_{\nu_g}u_f^g\,\dd S_g,
 \qquad f,h\in C_c^\infty(\Gamma),
\end{equation}
where $\nu_g$ is the outward unit normal. Thus the measured output
includes its boundary density. Equality on smooth supported data is
equivalent to equality of the maps: each element of
$H^{1/2}_{co}(\Gamma)$ has compact support in $\Gamma$ and can be
approximated in $H^{1/2}(\partial M)$ by smooth functions supported
in a fixed larger compact subset of $\Gamma$.

\begin{theorem}[Smooth anisotropic uniqueness from one patch]
\label{main:patch}
Let $n\ge3$. Let $M$ be a compact connected smooth $n$-manifold with
nonempty smooth boundary, and let $\Gamma\subset\partial M$ be nonempty,
proper, and relatively open. If two smooth Riemannian metrics $g_1,g_2$
on $M$ satisfy
\[
 \Lambda_{g_1,\Gamma}=\Lambda_{g_2,\Gamma},
\]
then there is a smooth diffeomorphism $\Phi:M\to M$, smooth through
the boundary, such that
\[
 \Phi|_\Gamma=\Id,\qquad g_2=\Phi^*g_1.
\]
\end{theorem}

Equality of the graphs of these operators is equivalent to the hypothesis.
Theorem~\ref{main:patch} resolves positively the smooth anisotropic
Calder\'on uniqueness problem with arbitrary same-patch measurements.
The datum is the density-valued energy map defined above.
The boundary may have several components and the manifold may be
nonorientable. The conclusion fixes every measured boundary point;
it imposes no pointwise condition on the inaccessible boundary.

The diffeomorphism ambiguity is unavoidable. If $\Phi$ fixes $\Gamma$
and $g_2=\Phi^*g_1$, then for $f$ supported in $\Gamma$ the pullback
$u_f^{g_1}\circ\Phi$ has boundary value $f$. Indeed a boundary
diffeomorphism fixing $\Gamma$ also preserves its complement.
Change of variables in the energy proves equality of the local maps.

\begin{corollary}[Full-boundary uniqueness]\label{main:full}
Let $M,g_1,g_2$ satisfy the geometric hypotheses of
Theorem~\ref{main:patch}. If their Dirichlet energies agree for every
pair of smooth boundary values, there is a smooth diffeomorphism
$F:M\to M$ fixing $\partial M$ pointwise and satisfying $g_1=F^*g_2$.
\end{corollary}

\begin{proof}
Fix $p\in\partial M$ and apply Theorem~\ref{main:patch} to
$\Gamma=\partial M\setminus\{p\}$. This is a nonempty proper open
patch, dense in the entire boundary. Continuity makes the resulting
$\Phi$ the identity on $\partial M$. Take $F=\Phi^{-1}$.
\end{proof}

There is also a consequence with no coordinate ambiguity.  Let
$\Omega\subset\mathbb R^n$ be bounded and connected with smooth
boundary, let $\Gamma\subset\partial\Omega$ be nonempty and relatively
open, and let $\gamma\in C^\infty(\overline\Omega)$ be strictly
positive.  For $f\in C_c^\infty(\Gamma)$, extended by zero to the
rest of the boundary, solve
\[
 \operatorname{div}(\gamma\nabla u_f^\gamma)=0\quad\text{in }\Omega,
 \qquad u_f^\gamma|_{\partial\Omega}=f.
\]
With $\nu_e$ the outward Euclidean unit normal, define the local
conductivity map by
\[
 \Lambda_\gamma^\Gamma f
       =\left.\gamma\partial_{\nu_e}u_f^\gamma\right|_\Gamma.
\]

\begin{corollary}[Smooth scalar uniqueness from one patch]
\label{main:scalar}
Let $n\ge3$, and let $\Omega$ and $\Gamma$ be as above.  If two
strictly positive conductivities
$\gamma_1,\gamma_2\in C^\infty(\overline\Omega)$ satisfy
\[
 \Lambda_{\gamma_1}^\Gamma f=\Lambda_{\gamma_2}^\Gamma f
       \quad\text{for every }f\in C_c^\infty(\Gamma),
\]
then $\gamma_1=\gamma_2$ on $\overline\Omega$.
\end{corollary}

The reduction in Section~\ref{scalar:section} identifies the exact
metric energy associated with a scalar conductivity.  It then proves
that a conformal self-diffeomorphism of the domain fixing a boundary
patch must be the identity.  No conductivity values on the inaccessible
boundary are assumed.

\subsection{Historical context}

Calder\'on formulated the inverse conductivity problem through boundary
measurements of Dirichlet energy and studied its linearization
\cite{Calderon1980}. Kohn and Vogelius established scalar boundary
determination \cite{KohnVogelius1984}. Sylvester and Uhlmann proved
full-data uniqueness for smooth scalar conductivities in Euclidean
dimensions at least three using complex geometrical optics solutions
\cite{SylvesterUhlmann1987}. For anisotropic conductivities, Lee and
Uhlmann developed boundary determination and proved real-analytic
recovery under geometric and topological hypotheses
\cite{LeeUhlmann1989}. Lassas and Uhlmann subsequently recovered a
real-analytic manifold and metric from an arbitrary nonempty
real-analytic boundary patch in dimensions at least three, without
the earlier topological restrictions
\cite[Theorem~1.1(ii)]{LassasUhlmann2001}. Their Green-function method
also led to the complete-manifold extension of Lassas, Taylor, and
Uhlmann \cite{LassasTaylorUhlmann2003}. Guillarmou and S\'a Barreto
proved same-patch uniqueness for smooth compact Einstein manifolds
with the same Einstein constant, recovering a boundary neighborhood
from the Einstein equation before applying analytic continuation in
the interior \cite[Theorem~1.1]{GuillarmouSaBarreto2009}.

For smooth scalar partial data, Bukhgeim and Uhlmann introduced a
Carleman-based recovery method \cite{BukhgeimUhlmann2002}. Kenig,
Sj\"ostrand, and Uhlmann prescribed input near a back face and
observation near a front face defined by a point outside the convex hull
\cite[Theorem~1.1]{KenigSjostrandUhlmann2007}; their conductivity
corollary also assumes equality on the whole boundary
\cite[Corollary~1.4]{KenigSjostrandUhlmann2007}. The two
measurement regions are constrained by this geometry. In dimension three,
Isakov obtained same-patch uniqueness by reflection when the inaccessible boundary lies in a plane
or sphere \cite{Isakov2007}. Dos Santos Ferreira, Kenig, Sj\"ostrand,
and Uhlmann proved the density-of-products theorem for the linearized
Schr\"odinger same-patch problem at zero potential
\cite[Theorem~1.1]{DosSantosFerreiraEtAl2009}. Alessandrini and Kim obtained
same-patch stability for scalar conductivities known near the whole
boundary \cite{AlessandriniKim2012}.

Smooth geometric results also exploit special backgrounds.
Dos Santos Ferreira, Kenig, Salo, and Uhlmann developed limiting
Carleman weights and recovered potentials and conformal factors on
admissible manifolds, which are conformal to submanifolds of a product
with a simple transverse factor \cite{DSFKSU2009}. Kenig and Salo
extended partial-data methods on such product geometries and related
them to geodesic and broken-ray transforms \cite{KenigSalo2013}.

Recent full-data results establish rigidity near prescribed smooth
backgrounds. Lin proves rigidity near the Euclidean metric under
smallness in a H\"older norm \cite[Theorem~1.1]{Lin2026}.
Mu\~noz-Thon and Stefanov prove local rigidity near every smooth
conformally Euclidean metric under Sobolev smallness
\cite[Theorem~2]{Stefanov2026}. Chen, Jiang, Liu, and Tao obtain
same-patch results under quasianalyticity, and smooth results with
additional foliation and one-sided ordering assumptions
\cite[Corollary~1.2, Theorem~1.6 and Corollary~1.7]{ChenJiangLiuTao2026}.
Theorem~\ref{main:patch} allows arbitrary smooth metrics and an
arbitrary common input and observation patch.

Uhlmann and Wang recover compactly supported potentials, and conformal
factors equal to one near the boundary, on a fixed near-Euclidean
three-dimensional metric with strictly convex boundary and an invertible
Dirichlet problem \cite{UhlmannWangNear2026}. Their partial-data result
recovers a potential near the measured boundary in a strictly convex
three-dimensional domain, with the potential supported away from the
measured patch \cite{UhlmannWangPartial2026}.
Daud\'e, Enciso, Helffer, Kamran, and Nicoleau propagate local DN
equality to a whole connected boundary component when the metrics
already agree in a collar of that component
\cite{DaudeEtAlPropagation2026}. These results clarify the separate
roles of a prescribed background, an initial region of coefficient
agreement, and the region reached by the measurements.

Harmonic functions provide an intrinsic way to describe points and
coordinates. Greene and Wu studied embeddings by harmonic functions
\cite{GreeneWu1975}; the Poisson embedding approach of Lassas,
Liimatainen, and Salo represents a point by its evaluation on
boundary-generated harmonic functions \cite{LLS20}. Their source
embedding treats solutions generated in a remote interior region
\cite[Section~6.1]{LLS20}. The smooth separation, approximation and
finite-jet arguments are local tools; their global continuation results
in dimensions at least three require real analyticity. Here the smooth
continuation step is supplied by
Theorem~\ref{local:extension}; its uniform estimate on shrinking
surfaces is the central analytic argument.

The approximation mechanism goes back to Lax and Malgrange
\cite{Lax1956,Malgrange1956}. R\"uland and Salo developed its
quantitative supported-boundary form \cite{RulandSalo2019}.
We use point-distribution duality to obtain just the finite harmonic
jets needed for the local argument. Localized boundary determination
and exterior data transfer have further precedents in
\cite{AlessandriniGaburro2009,AlessandriniKim2012}.

\subsection{Proof strategy}

The global argument is organized around the local metric extension
Theorem~\ref{local:extension}. Its input is a pair of interior harmonic
coordinate charts in which the metrics, their Green kernels, and a
finite family of paired harmonic functions already agree on one side of
a smooth hypersurface. The first-side Hessians span all second derivatives
allowed by the harmonic equation. The theorem extends metric equality
across the hypersurface. The main task is to prove this continuation theorem for
smooth coefficients.

First, the measured energy form supplies the interior data needed to
start. Local boundary determination recovers the full boundary jets
and permits an exterior cap on which the metrics agree. A variational
comparison then makes the actual Dirichlet Green kernels agree on the
cap. Harmonic functions generated by sources in a remote part of that
cap separate interior points and supply all admissible second-order
harmonic jets. They will identify the paired charts and the finite
harmonic family required by the local theorem.

Inside these charts, a mixed Green kernel solves the first metric
equation in its first variable and the second metric equation in its
second variable. Product continuation constructs this kernel at fixed
positive separation from the diagonal. Common-basis methods for this
separately harmonic extension have antecedents in \cite{Zeriahi1982}.
Its flux on a small sphere
centered at $y$ defines a functional $T_y$ that maps local first-metric
harmonic functions to local second-metric harmonic functions. It sends
each chosen first-side function to its paired second-side function and
reproduces constants and coordinates: $T_y1=1$ and $T_yx^i=y^i$.
Existence at each positive radius is not yet a bound as that radius
tends to zero.

To compare conformal classes, we normalize the inverse metrics to have
the same trace against a reference metric conformal to the first.
The shrinking argument temporarily assumes that their difference,
together with a fixed finite number of its derivatives, is small
compared with the sphere radius $r$.
Under this assumption, we construct the transfer functional and bound
its representing density uniformly in $L^2$ over a fixed set of
centers and their unit spheres. Exact reproduction of constants and coordinates cancels the
affine Taylor polynomial of a harmonic function. Its remainder has
size $O(r^2)$ on the sphere, so for a paired harmonic difference $V$
we obtain an $L^2$ bound for $V/r^2$. Finite-order interpolation then
bounds the needed derivatives of $V$ by $Cr^{3/2}$. The spanning
harmonic Hessians turn this into the same superlinear bound for the
normalized metric discrepancy and its required derivatives. It is a
strict improvement of the temporary bound at small radii.

There are two scales in this argument. A spatial dilation is chosen
once to make the initial metric discrepancy small and to handle radii
above a fixed threshold. The strict improvement then allows the radius
parameter to tend to zero with that dilation held fixed. The radii
vanish on a fixed neighborhood, forcing equality of the normalized
inverse metrics there. This determines the metrics up to a conformal
factor. The common harmonic coordinates force that factor to be
constant when $n\ge3$, and agreement on the known side fixes it.

The uniform density estimate must control the full equation obtained
by differentiating traces on the moving spheres. Comparing spherical
harmonic degree-raising and degree-lowering operators gives a positive
weighted estimate, using the tensor and spherical-harmonic framework
that also appears in tensor tomography \cite{PSU15,GPSU16}.
The moving geometry contributes second-order
terms, which cannot simply be treated as lower-order errors. A second
comparison retains those terms and bounds their contribution by the
positive quantities in the first estimate.

Finally, the source evaluations make the local identifications
single-valued and prevent an interior point from escaping to the
boundary. A ball in the matched region touching its frontier supplies
the one-sided geometry of the local theorem. Applying that theorem
eliminates every interior frontier. The resulting isometry extends
smoothly through the boundary, and the original supported boundary
data force it to fix all of $\Gamma$.

Section~\ref{data:section} constructs the common cap and harmonic tests.
Section~\ref{cross:section} states the local extension theorem and
constructs product continuations at fixed separation;
Section~\ref{sec:balls} uses them to construct the sphere functionals.
Section~\ref{sec:angular} proves the coercive estimate and bounds the
transfer density. Section~\ref{sec:bootstrap} turns the resulting
moment bound into local metric extension.
Section~\ref{global:section} gives the global identification and the
boundary conclusion. Section~\ref{scalar:section} proves the scalar
conductivity consequence.
